\documentclass[conference]{IEEEtran}
\IEEEoverridecommandlockouts

\usepackage{amsmath,amssymb,amsfonts}
\usepackage{amsthm}
\usepackage{graphicx}
\usepackage{textcomp}
\usepackage{xcolor}
\usepackage{booktabs}
\usepackage{tabularx}
\usepackage{makecell}
\usepackage{adjustbox}

\usepackage[style=numeric, backend=biber]{biblatex}
\addbibresource{references.bib}

\usepackage{algorithm}
\usepackage{algpseudocode}

\theoremstyle{definition}
\newtheorem{definition}{Definition}

\theoremstyle{plain}
\newtheorem{theorem}{Theorem}
\newtheorem{lemma}{Lemma}
\newtheorem{proposition}{Proposition}

\usepackage{microtype}
\usepackage{listings}

\def\BibTeX{{\rm B\kern-.05em{\sc i\kern-.025em b}\kern-.08em
    T\kern-.1667em\lower.7ex\hbox{E}\kern-.125emX}}

% Configuration for code listings
\definecolor{codegray}{rgb}{0.5,0.5,0.5}
\definecolor{codeback}{rgb}{0.97,0.97,0.97}
\lstset{
    backgroundcolor=\color{codeback},
    basicstyle=\ttfamily\small,
    breaklines=true,
    captionpos=b,
    commentstyle=\color{codegray},
    frame=single,
    keepspaces=true,
    numbers=none,
    keywordstyle=\color{black}\bfseries,
    showspaces=false,
    showstringspaces=false,
    showtabs=false,
    tabsize=4
}

\usepackage{hyperref}

\begin{document}

\title{pending}

\author{
\IEEEauthorblockN{Chaolu Liu}
\IEEEauthorblockA{\textit{School of Computer Science} \\
\textit{University College Dublin}\\
Dublin, Ireland \\
cl5248@hotmail.com}
\and
\IEEEauthorblockN{Emilio Martinez}
\IEEEauthorblockA{\textit{School of Computing} \\
\textit{National College of Ireland}\\
Dublin, Ireland \\
eyolticmg@icloud.com}
\and
\IEEEauthorblockN{Furqan Rustam}
\IEEEauthorblockA{\textit{School of Computer Science} \\
\textit{University College Dublin}\\
Dublin, Ireland \\
furqan.rustam1@gmail.com}
\and
\IEEEauthorblockN{Anca Delia Jurcut}
\IEEEauthorblockA{\textit{School of Computer Science} \\
\textit{University College Dublin}\\
Dublin, Ireland \\
anca.jurcut@ucd.ie}
}

\maketitle

\begin{abstract}
pending
\end{abstract}
\begin{IEEEkeywords}
Poseidon, Poseidon2, temporal diffusion, algebraic reachability,
subspace trails, algebraic cryptanalysis, zero-knowledge proofs
\end{IEEEkeywords}

\section{Introduction}
\label{sec:introduction}

\subsection{Problem Statement}
Given a fully specified permutation, an application interface, and a
resource budget, determine whether a candidate low-growth trajectory
can be reached from admissible inputs, traversed with a checkable
algebraic certificate, and used to meet the required output constraints.
The output should separate a structural witness from an attack witness
and compare total attack cost with an appropriate baseline.

The analysis has three objectives: characterize round-dependent inactive
directions; determine constrained entry and exit feasibility; and assess
whether exploiting a certified trajectory improves the specified attack
task. Structural exclusions remain scoped to the tested model.
\textcolor{red}{pending text: Establish a distinct technical improvement
for a declared class of instances and complete the corresponding
fixed-instance attack or exclusion argument.}

\textcolor{blue}{How to reconcile cryptographic safety against advanced algebraic attacks, such as subspace trails and full-round absorption, with the lightweight arithmetization efficiency in ZK-friendly permutations?}


\section{Background}
\label{sec:background}
Let $p$ be prime and $x_r\in\mathbb F_p^t$. For the analysis below use
\begin{equation}
x_{r+1}=R_r(x_r)=M_rS_r(x_r+c_r)+d_r,
\label{eq:round}
\end{equation}
where $M_r\in\mathrm{GL}_t(\mathbb F_p)$ and $c_r,d_r\in\mathbb F_p^t$.
For an active set $A_r\subseteq\{0,\ldots,t-1\}$, let
\[
(S_r(z))_j=
\begin{cases}
z_j^{\alpha_r},&j\in A_r,\\
z_j,&j\notin A_r,
\end{cases}
\quad \gcd(\alpha_r,p-1)=1.
\]
Full rounds activate every coordinate; partial rounds activate a proper
subset. The source convention $M_r(S_r(x_r)+C_r)$ is recovered by
$c_r=0$, $d_r=M_rC_r$.
Actual instances must preserve their specified operation order and
initial mixing. These power-map arguments do not automatically apply to
different nonlinear layers.
\textcolor{red}{pending text: Provide tested specification adapters,
version hashes, and known-answer vectors for each evaluated instance,
including Neptune if retained.}

Let $E_r$ select the active coordinates and $W_r=\ker E_r$. Write
$F_{a,b}=R_{b-1}\circ\cdots\circ R_a$ and $F_{a,a}=I$.
An affine family $a+Bz$ uses a full-column-rank basis
$B\in\mathbb F_p^{t\times k}$ and $z\in\mathbb F_p^k$.

For degree and diffusion. The branch number is
\[
B_d(M)=\min_{v\ne0}\bigl(\mathrm{wt}(v)+\mathrm{wt}(Mv)\bigr).
\]
An MDS matrix has branch number $t+1$. This one-step support property
does not exclude multiround cancellations or constrained directions.
Distinguish formal polynomial degree, degree after family substitution,
and degree after field-equation reduction. Degree-propagation recurrences
are upper bounds unless tightness is proved; cancellation can occur.
These quantities alone do not determine the solving degree of a
Gr\"obner-basis computation.

\section{Related Work}
\label{sec:related}

The security analysis of ZK-friendly permutations heavily relies on characterizing the propagation of algebraic properties through rounds with partial non-linear layers. Early investigations primarily focused on identifying invariant spaces and establishing sufficient conditions to prevent infinitely long subspace trails under stationary linear layers~\cite{weaklayers2020, algebraiccrypt2023}. Complementing these structural checks, Keller and Rosemarin re-evaluated the security margins of the HADES design strategy by examining the middle layers~\cite{mindthemiddle2021}. More recently, research has shifted toward finite-horizon and dynamic vulnerabilities. Li and Wang formalized block-matrix criteria for finite linear subspace trails, demonstrating their unavoidability in standard designs~\cite{cryptoeprint:2026/967}. Concurrently, Grassi, Koschatko, and Rechberger explored trail-assisted Gröbner basis solving strategies across various operational modes~\cite{poseidonneptune2025}. Regarding the reachability of these internal trajectories through full-round prefixes, Vitto introduced the concept of full-round absorption and adaptive weak matrix construction (Slipway), showing how constrained families can bypass initial degree growth~\cite{slipway2026}. Furthermore, Li, Liu, and Wang extended these structural insights beyond linear boundaries by developing non-linear subspace methods for advanced algebraic attacks~\cite{nonlinear2026}. While existing literature thoroughly characterizes individual trail structures and specific attack vectors, our work builds upon these frameworks by formalizing temporal diffusion networks and examining the rigorous algebraic reachability conditions governing constrained input-output mapping.


\section{Threat Model}
\label{sec:model}

\subsection{Threat Model and Attack Tasks}
For a fixed-instance attack, the public matrix, constants, and mode
are fixed before target selection. The adversary may not replace them.
Given admissible inputs $\mathcal I$ and required outputs $\mathcal O$,
the constrained-input, constrained-output (CICO) task is to find
\begin{equation}
x\in\mathcal I,\qquad F_{0,R}(x)\in\mathcal O.
\label{eq:cico}
\end{equation}
Record fixed coordinates, target distribution, and whether targets are
adversarially chosen. Hash preimages and collisions must additionally
respect the message interface, padding, domain separation, rate/capacity,
and any compression feed-forward. CICO is not an automatic reduction
to every mode~\cite{poseidonneptune2025,cryptoeprint:2026/306}.
\textcolor{red}{pending text: Instantiate attack games, generic comparison
algorithms, success criteria, and concrete resource bounds. A protocol
forgery claim requires a specified protocol and a verifier-accepting
transcript violating the stated security goal.}
\subsection{Success Criteria and Cost Accounting}
Restricted degree alone does not imply proportional reduction in solving
degree or practical attack cost. Variables, sparsity, term order, field
equations, and solver implementation also matter. Record
\begin{equation}
\begin{split}
C_{\rm total}={}&C_{\rm search}+C_{\rm entry}+C_{\rm solve}\\
&+C_{\rm output}+C_{\rm verify},
\end{split}
\label{eq:cost}
\end{equation}
with repetition counts, success probability, memory, and amortized
precomputation separately stated.
\textcolor{red}{pending text: Establish a justified complexity model and
measured end-to-end advantage. The source's algebraic-complexity-collapse
theorem is not established and is not retained as a result.}


\section{Methodology}
\label{sec:methodology}
\label{sec:theory}
The technical approach has three components. Temporal analysis produces
candidate direction spaces over a declared round interval. Reachability
analysis then tests compatibility with the input, prefix, suffix, and
output constraints. Certificate-producing analysis records the resulting
scoped claims and checks witnesses independently.
\subsection{Temporal Diffusion Networks}
\begin{definition}[Support TDN]
For interval $[a,b)$, create vertices $(r,j)$ at round boundaries.
Include $(r,j)\to(r+1,i)$ if $(M_r)_{ij}\ne0$, and label the
nonlinear operation at $(r,j)$ when $j\in A_r$.
\end{definition}
This graph represents possible dependencies in the coordinatewise round
model. Constants do not introduce dependence on family parameters.
Graph reachability overapproximates algebraic dependence: weighted paths
can cancel, and a direction can have zero projection on an active
coordinate despite existing paths.
The intended exact TDN layer stores transition matrices, selectors, and
surviving direction bases.
\textcolor{red}{pending text: Define and prove any stronger correspondence
between TDN properties and algebraic trajectories. A support graph alone
is not an exact certificate.}
\subsubsection{Finite-Horizon Inactive Directions}
Set $\Phi_{a,a}=I$ and $\Phi_{a,r}=M_{r-1}\cdots M_a$ for $r>a$.
Define
\begin{equation}
H_{a,b}=
\begin{bmatrix}
E_a\Phi_{a,a}\\ E_{a+1}\Phi_{a,a+1}\\ \vdots\\
E_{b-1}\Phi_{a,b-1}
\end{bmatrix},
\qquad K_{a,b}=\ker H_{a,b}.
\label{eq:observability}
\end{equation}
The following elementary identities are background linear algebra,
not new cryptanalytic contributions.

\begin{proposition}[Affine propagation]
\label{prop:affine}
If the columns of $B$ lie in $K_{a,b}$, then for every base state $q$,
\begin{equation}
F_{a,r}(q+Bz)=F_{a,r}(q)+\Phi_{a,r}Bz
\label{eq:affine}
\end{equation}
for all $z$ and $a\le r\le b$.
\end{proposition}
\begin{proof}
The identity holds at $r=a$. At each boundary $r<b$,
$E_r\Phi_{a,r}B=0$, so every family variation $v=\Phi_{a,r}Bz$
is zero at the active coordinates. Thus $S_r(y+v)=S_r(y)+v$
for every $y$. Substitution into~\eqref{eq:round} multiplies the
variation by $M_r$, while constants cancel. Induction proves the claim.
\end{proof}
The degree is at most one in the family parameters during this interval,
not necessarily in unrestricted inputs to the whole permutation.
The base state's active coordinates may be nonzero. No condition that
round constants lie in the direction space is needed.

\begin{proposition}[Dimension]
\label{prop:dimension}
If $H_{a,b}$ has $m$ rows, then
$\dim K_{a,b}=t-\operatorname{rank}(H_{a,b})\ge\max(0,t-m)$.
In particular, $h<t$ one-active-coordinate rounds always have a
nonzero certified direction, including with MDS mixing.
\end{proposition}
\begin{proof}
Apply rank--nullity and the bound
$\operatorname{rank}(H_{a,b})\le\min(t,m)$.
\end{proof}
Consequently, a short finite inactive-direction trail is not by itself
evidence of a defective matrix.

\subsubsection{Stationary Layers and Exact Controls}
For stationary $M_r=M$, $E_r=E$, set
\begin{equation}
H_t=[E^\top,(EM)^\top,\ldots,(EM^{t-1})^\top]^\top,
\quad K_\infty=\ker H_t.
\label{eq:stationary}
\end{equation}
\begin{proposition}[Largest inactive invariant subspace]
$K_\infty=\bigcap_{j\ge0}\ker(EM^j)$ is the largest $M$-invariant
subspace contained in $\ker E$.
\end{proposition}
\begin{proof}
Cayley--Hamilton expresses every higher power as a linear combination
of $I,M,\ldots,M^{t-1}$, giving the intersection identity.
If $v$ lies in that intersection, so does $Mv$. Conversely, any
$M$-invariant $V\subseteq\ker E$ satisfies $EM^jV=0$ for all $j$.
\end{proof}
This is a certificate for inactive directions in the stated model,
not an exclusion of active or nonlinear trails~\cite{weaklayers2020,nonlinear2026}.
For stationary one-coordinate activation define
$\tau(v)=\min\{j\ge0:EM^jv\ne0\}$, with value $\infty$ if empty.
If $\tau(v)=d$, the first $d$ rounds are affine along $v$.
Subsequent activation does not prove a prescribed global degree or cost.

Exact controls are available without experiments.
For $M=I$, $E=e_0^\top$, the kernel has dimension $t-1$.
For the cyclic shift
$M(x_0,\ldots,x_{t-1})^\top=(x_1,\ldots,x_{t-1},x_0)^\top$,
the first $t$ rows $EM^j$ select all coordinates.
Thus $K_\infty=0$ and the horizon-$h$ kernel has dimension $t-h$
for $h\le t$. These are mathematical controls, not measured results.
\subsection{Constrained Algebraic Reachability}
\label{sec:reachability}
For candidate affine sets $q_r+\operatorname{im}B_r$, individual-state
reachability is the feasibility problem
\begin{equation}
\begin{split}
x_0&\in\mathcal I,\quad x_R\in\mathcal O,\\
x_{r+1}&=R_r(x_r)\quad(0\le r<R),\\
x_r&=q_r+B_rz_r\quad(a\le r\le b).
\end{split}
\label{eq:reachability}
\end{equation}
One solution is weaker than a useful reachable family.
Family reachability additionally requires a parameterization $\psi(u)$
with $\psi(u)\in\mathcal I$ and
$F_{0,a}(\psi(u))=q_a+B_a\eta(u)$ on the stated parameter domain.
Report the image dimension of $\eta$, the representation degree, and
the cost of constructing or evaluating $\psi$. A high-degree prefix
can move the difficulty rather than remove it.

Permutation prefixes give every internal state a unique unrestricted
preimage. They do not guarantee an admissible preimage or a cheap family
description. Entry is a linear system only when the prefix and imposed
constraints are affine; nonlinear prefixes generally yield polynomial
systems.
\textcolor{red}{pending text: Prove new reachability or exclusion results,
or exhibit independently verified improvements beyond prior attacks.
Reachability itself cannot be claimed as new in view of Slipway.}
\subsubsection{Nonlinear Trajectory Certificates}
A proposed extension uses varieties $V_{\mathbb F_p}(I_r)$.
For every generator $g$ of $I_{r+1}$, membership
$g\circ R_r\in I_r$ suffices to certify
$R_r(V_{\mathbb F_p}(I_r))\subseteq V_{\mathbb F_p}(I_{r+1})$:
evaluate the ideal-membership identity on the source zero set.
Plain ideal membership need not characterize all finite-field inclusions
without appropriate vanishing ideals or field equations.
\textcolor{red}{pending text: Specify a tractable certificate class,
prove its scope, and compare with nonlinear-subspace methods.}
\subsection{Certificate-Producing Analysis}
\label{sec:framework}
\subsubsection{Finite-Horizon Scanner}
The proposed algorithm replaces the source scanner's binary security verdict.
\begin{algorithm}[t]
\small
\caption{Finite-horizon inactive-direction certificate}
\label{alg:subspace_scanner}
\begin{algorithmic}[1]
\Require Exact field, interval $[a,b)$, matrices $M_r$, selectors $E_r$
\State $\Phi\gets I_t$; $H\gets$ empty matrix with $t$ columns
\For{$r=a,\ldots,b-1$}
  \State Append rows $E_r\Phi$ to $H$
  \State $\Phi\gets M_r\Phi$
\EndFor
\State $B\gets\operatorname{KernelBasis}(H)$ over $\mathbb F_p$
\State Verify $HB=0$ and $\operatorname{rank}B=t-\operatorname{rank}H$
\State \Return $(B,\operatorname{rank}H,\text{interval},\text{model hash})$
\end{algorithmic}
\end{algorithm}
This is an algorithm specification, not a claim of completed software.
For stationary infinite horizons use~\eqref{eq:stationary}.
Outputs mean a certified nonzero inactive-direction space, zero such
space in the specified model, or inconclusive/error.
Zero kernel is not a general \texttt{SECURE} status.

The lower-right block of $M$ need not be a restriction to $\ker E$,
because that space need not be invariant. Characteristic polynomial
factorization may remain diagnostic; kernels of individual factors are
not assumed to enumerate every relevant invariant subspace.
Constants still matter for the base trajectory and constrained entry.
\subsubsection{Witness Checking and Completeness}
Build~\eqref{eq:reachability} with explicit intermediate-state equations.
Record the selected elimination or solving method.
Independently evaluate every solution using the original permutation
and mode interface. Family identities require symbolic certificates or
proofs; sampling alone is not a proof. Exclusion claims need a complete
argument or checkable unsatisfiability certificate.
Timeout means inconclusive.
\textcolor{red}{pending text: Implement the checker and independent
verifier; release accepted and rejected certificates, known-answer tests,
and verified attack witnesses.}
\subsubsection{Complexity and Scalability}
For horizon $h$, classical dense transition products cost $O(ht^3)$
field operations. If $m=\sum_r|A_r|$, a conservative elimination bound
for the $m\times t$ stack is $O(mt^2+t^3)$, with
$O(mt+t^2)$ field-element storage before certificate overhead.
For one-coordinate activation, $m=h$. These are bounds for a
straightforward implementation; bit cost depends on the modulus.
\textcolor{red}{pending text: Measure runtime, memory, certificate size,
and verification overhead. No sub-millisecond timing or negligible
integration overhead is claimed.}


\section{Adversarial Hash Design}
\label{sec:adversarial}
Adversarial Hash Design supplies controlled test instances for the
methodology in Section~\ref{sec:methodology}. Its role is to challenge
candidate structures and study cost/attack trade-offs; it is not a
separate claim that an optimized candidate is secure.
\subsection{Candidate Generation and Held-Out Evaluation}
Generate layers or schedules within declared invertibility and cost
constraints. Separate search/training instances from held-out targets
and seeds; archive all candidates and selection decisions.
\subsection{Design Freedom and Fixed-Instance Attacks}
In a fixed-primitive attack, matrices and constants are fixed in advance.
In a design-space experiment, disclose whether constants and targets were
known when selecting a matrix. A constant-dependent matrix construction
is not a break of a fixed published instance.
\subsection{Screening Objectives and Selection}
Search objectives may include field-operation counts, certified inactive
dimensions, entry cost, and verified solving cost. Passing a screen
establishes only the screened property.
\textcolor{red}{pending text: Provide search costs, held-out performance,
overfitting checks, and a verified Pareto comparison. A concrete new
permutation is optional and is not recommended for deployment here.}

Every selected candidate should retain the exact matrices, constants,
schedule, selection seed, search budget, and certificates produced by
Section~\ref{sec:framework}. Failed entry searches, inconclusive solver
runs, and rejected witnesses must remain in the experimental record.
Section~\ref{sec:evaluation} evaluates selected candidates on the same
tasks and budgets as the declared controls.


\section{Evaluation}
\label{sec:evaluation}
\subsection{Research Questions and Instances}
Test whether exact temporal analysis improves on support-only metrics;
which trails admit constrained entry and exit; whether guided solving
improves total cost; and how design cost relates to measured attack cost.
\textcolor{red}{pending text: Answer these questions with certificates
and matched experiments, including null outcomes.}

Use exhaustive small-field controls, solver-feasible intermediate
instances, and exact structural checks over production fields.
The proposed width grid is $t\in\{3,4,8,12,16,24,32\}$ where supported
by each selected specification; include the BN254 scalar field from
the source draft.
\textcolor{red}{pending text: Freeze exact primes, valid exponents,
round schedules, constants, versions, and modes in a manifest.}
Small-field attack costs must not be treated as production-field costs.

Include published instances, dense/MDS controls, identity, cyclic-shift,
block, planted-defect, random, and adversarially selected families.
Distinguish genuine published matrices from superficially similar replacements.
A proposed corpus budget is $10^4$ candidates per selected family/width
cell, subject to pilot feasibility.
\textcolor{red}{pending text: Fix sample counts, distributions, coefficient
domains, seeds, rejection rules, and duplicate policy. Report all candidates
and invertibility rejections; no corpus has been established by the supplied data.}

\subsection{Baselines and Ablations}
Compare direct algebraic modeling, support propagation, stationary
invariant tests, finite linear trails, and applicable published attacks
\cite{poseidonneptune2025,cryptoeprint:2026/306,slipway2026,nonlinear2026}.
Match fields, round schedules, tasks, target distributions, and limits.
Ablate changing matrices, activation schedules, entry/output constraints,
and nonlinear restrictions separately. Match arithmetic cost or S-box
budget where useful, disclosing unavoidable trade-offs.
\textcolor{red}{pending text: Reproduce the closest baselines and provide
paired success, time, and memory comparisons.}

\subsection{Measurements and Results}
Record kernel dimensions and ranks, interval lengths, reachable-family
dimensions, entry/output costs, success probabilities, total wall time,
CPU/GPU time, peak memory, and certificate size. Distinguish exact degree
from upper bounds and specify variables and reductions.
Treat timeouts and memory failures as censored outcomes.

\begin{table}[t]
\centering
\caption{Planned structural and reachability evaluation. No empirical
entries are available yet.}
\label{tab:summary_all}
\small
\begin{tabularx}{\columnwidth}{@{}lX@{}}
\toprule
Family & Required evidence\\
\midrule
Published & \textcolor{red}{pending text: Instance IDs; kernel profiles;
reachable families}\\
Dense/MDS & \textcolor{red}{pending text: Matched controls and certificates}\\
Sparse/block & \textcolor{red}{pending text: Planted and unplanted outcomes}\\
Adversarial & \textcolor{red}{pending text: Held-out outcomes; search costs}\\
\bottomrule
\end{tabularx}
\end{table}
\begin{table}[t]
\centering
\caption{Planned end-to-end comparisons. No speedups are asserted.}
\small
\begin{tabularx}{\columnwidth}{@{}lX@{}}
\toprule
Comparison & Required evidence\\
\midrule
Direct vs.\ guided & \textcolor{red}{pending text: Success, time, memory}\\
Prior vs.\ proposed & \textcolor{red}{pending text: Matched tasks and limits}\\
Entry/exit & \textcolor{red}{pending text: Separate and total costs}\\
Verification & \textcolor{red}{pending text: Witnesses and checking time}\\
\bottomrule
\end{tabularx}
\end{table}

\subsection{Resources and Initial Three-Month Campaign}
The discussion describes approximately EUR~30,000 of workstation
resources and an accelerator with approximately 128~GB memory.
These are planning inputs, not verified specifications.
\textcolor{red}{pending text: Record CPU, host RAM, GPU model/memory,
storage, software versions, arithmetic backends, and compiler settings.}
Profile symbolic workloads before assuming GPU acceleration.

A proposed twelve-week initial campaign is:
\begin{enumerate}
\item Weeks 1--2: freeze definitions, validate exact controls, reproduce
reference instances, and pilot resource use.
\item Weeks 3--4: complete scoped certificates, independent checking,
and baseline reproduction.
\item Weeks 5--8: run preregistered corpora, entry/exit studies, and paired solves.
\item Weeks 9--10: run held-out adversarial candidates and ablations.
\item Weeks 11--12: verify witnesses, analyze uncertainty, and package artifacts.
\end{enumerate}
This schedule does not guarantee a major-conference contribution.
\textcolor{red}{pending text: Freeze per-job limits, concurrency,
checkpointing, stopping rules, and actual run logs.}

\subsection{Reproducibility}
Release manifests, parameter files, seeds, versions, raw logs, failures,
and a witness checker. Separate discovery from independent verification,
record manual interventions, and report negative results.
For random corpora report uncertainty conditional on the sampling model;
optimized samples do not estimate random-population prevalence.
\textcolor{red}{pending text: Provide repository/archive identifiers,
license, checksums, reproduction instructions, and an independent rerun.
No public tool release is claimed.}


\section{Discussion \& Limitations}
\label{sec:discussion}
\subsection{Meaning and Limits of Certificates}
An internal affine family may reduce dependence on some variables without
making admissible inputs easier to find. A zero kernel excludes only
the specified inactive-direction model. Branch number matters but cannot
replace multiround, mode-aware analysis; finite trails alone are not
enough to reject a design.
\subsection{Protocol and Production-Parameter Limits}
State whether results concern reduced rounds, chosen matrices, fixed
published parameters, or a specific mode. Small-instance successes
\cite{cryptoeprint:2026/150} do not themselves extrapolate to production
attacks. A permutation weakness does not establish proof forgery.
The cited Fiat--Shamir study reports negative results for its explored
approaches~\cite{cryptoeprint:2026/409}, so it cannot support an automatic
forgery conclusion.
\textcolor{red}{pending text: Establish any retained protocol consequence
with a complete attack; otherwise restrict conclusions to verified
permutation/hash tasks.}
\subsection{Unresolved Evidence and Scope}
The proposed evaluation has not supplied measured timings, failure
rates, validated attack advantages, or completed tool artifacts.
The linear-algebra propositions concern the stated partial power-map
model and do not cover all nonlinear trajectory classes.
\textcolor{red}{pending text: Complete theorem-level positioning,
implementation validation, matched measurements, and independent
witness checks before making broader claims. Confirm author details
and venue anonymization requirements before submission.}


\section{Conclusion}
pending

\printbibliography


\appendices
\section{Notes}
\textcolor{red}{pending text: Confirm author details and anonymization
requirements; review proofs; finish related work; freeze manifests;
complete results; archive artifacts; rewrite abstract and conclusion.}


\end{document}
